\section{Comparison and the tracial ultrapower}
\label{u:sec:preliminaries}

We now begin the direct unital comparison construction. Its input is the strict comparison in Theorem~\ref{u:thm:main}; its output will be uniform property $\Gamma$ without a real-rank-zero assumption. This section fixes the distinct ultrapower $D=A^\omega$, its limit traces, and the bounded diagonal principle used to turn approximate coefficient tests into exact identities.

For a $C^*$-algebra $E$, write $E_+$ for its positive cone, $E_{\mathrm{sa}}$
for its selfadjoint part, and
\[
 \Her(b)=\overline{bEb}\qquad(b\in E_+)
\]
for the hereditary subalgebra generated by $b$. We use
$(b-t)_+$ for the positive part of $b-t1$, computed in the unitization
when necessary. Spectral projections, including
$\supp(b)=1_{(0,\infty)}(b)$, are taken in the indicated bidual or von
Neumann algebra. All tensor products in the proof are minimal.

\subsection{Traces and comparison conventions}

Let $\K$ be the compact operators on a separable infinite-dimensional
Hilbert space. For $a,b\in(A\otimes_{\min}\K)_+$, Cuntz subequivalence
$a\precsim b$ means that there are $x_n\in A\otimes_{\min}\K$ with
\[
 \norm{x_n^*bx_n-a}\longrightarrow0.
\]
Cuntz equivalence is mutual subequivalence. Its classes form the
ordered abelian semigroup $\Cu(A)$, with $[a]\le[b]$ when $a\precsim b$
and addition given by orthogonal direct sum.

The tracial-state space is denoted by $T(A)$. Each $\tau\in T(A)$
extends to the usual lower semicontinuous trace $\tau\otimes\Tr$ on
$A\otimes_{\min}\K$, where $\Tr$ is unnormalized. Set
\[
 d_\tau(a)=\lim_{n\to\infty}(\tau\otimes\Tr)(a^{1/n})
       \in[0,\infty].
\]
In particular, on $M_j(A)$ the trace of the unit is $j$, not $1$.
Matrix traces will have this normalization throughout.

\begin{definition}[Strict comparison]
\label{u:def:strict-comparison}
The algebra $A$ has strict comparison if, whenever
$a,b\in(A\otimes_{\min}\K)_+$ and $b\ne0$, the inequalities
\[
 d_\tau(a)<d_\tau(b)\qquad(\tau\in T(A))
\]
imply $a\precsim b$. When $T(A)=\varnothing$, this requires
$a\precsim b$ for every positive $a$ and every nonzero positive $b$
in the stabilization.
\end{definition}

Suppose for now that $A$ is simple and unital and that
$T:=T(A)\ne\varnothing$. The set $T$ is weak star compact: it is the
closed subset of the state space defined by the trace identities.
Every $\tau\in T$ is faithful. Indeed,
$\{x\in A:\tau(x^*x)=0\}$ is a closed two-sided ideal, by the trace
identity and the inequalities for multiplication in the tracial
$2$-seminorm, and it does not contain $1$.

We specify the trace extensions used in spectral calculations. The
bounded trace $\tau\otimes\Tr_j$ on $M_j(A)$ has its canonical normal
extension to $M_j(A)^{**}$. It is tracial, as follows by extending the
trace identity in each variable using separate weak star continuity.
We use the same symbol for this extension. Spectral calculus gives
\[
 d_\tau(v)=\tau(\supp(v))\qquad(v\in M_j(A)_+).
\]
For the entire stabilization one may instead work in
$M_\tau\,\overline\otimes\,\mathcal B(\ell^2)$, where
$M_\tau=\pi_\tau(A)''$ is the tracial GNS von Neumann algebra, with
its normal semifinite trace induced by $\tau\otimes\Tr$. The same
support formula holds there, with the value $\infty$ permitted. To
see it, first rescale a positive element to a contraction and then
use monotone convergence on its spectral resolution. Scalar
rescaling leaves the limiting rank unchanged.

Here is the comparison convention needed for the known absorption
criterion. A functional on $\Cu(A)$ is an additive, order-preserving
map $\varphi:\Cu(A)\to[0,\infty]$ that vanishes at $0$ and preserves
suprema of increasing sequences; let $F(\Cu(A))$ denote these
functionals. The functional formulation of strict comparison used
in \cite[Section~5, pp.~20--21 of the arXiv version]{CETW} requires
$[u]\le[v]$ if
\[
 \varphi([u])\le\varphi([v])\quad\text{for every }
 \varphi\in F(\Cu(A)),
 \qquad
 \varphi([u])<\varphi([v])\quad\text{whenever }
 0<\varphi([v])<\infty.
\]

\begin{lemma}[Compatibility of comparison conventions]
\label{u:lem:comparison-convention}
Let $A$ be simple and unital with $T(A)\ne\varnothing$. Strict
comparison in Definition~\ref{u:def:strict-comparison} implies the
functional formulation just stated.
\end{lemma}

\begin{proof}
For every $\tau\in T$, the map $[a]\mapsto d_\tau(a)$ is a
functional on $\Cu(A)$, by \cite[Proposition~4.2]{ERS}, applied to
the lower semicontinuous trace $\tau\otimes\Tr$. Only this direction
of the trace--functional correspondence is needed.

Let $u,v\in(A\otimes_{\min}\K)_+$ satisfy the functional
inequalities. If $v=0$, any $\tau\in T$ gives $d_\tau(u)=0$.
The faithful stabilized trace then gives $u=0$, so $u\precsim v$.
Assume $v\ne0$. Faithfulness implies $d_\tau(v)>0$ for every
$\tau\in T$. Thus
\[
 d_\tau(u)<d_\tau(v)
 \quad\text{whenever }d_\tau(v)<\infty.
\]

Fix $s>0$. There is a positive element $c\in M_k(A)$ in a finite
matrix corner such that $\norm{u-c}<s/2$: compress $u$ by a
sufficiently large finite-rank corner projection. In the tracial
von Neumann representation for any $\tau\in T$, put
$P=1_{(s,\infty)}(u)$. Then
\[
 PcP\ge\frac{s}{2}P.
\]
The polar decomposition of $c^{1/2}P$ therefore has initial
projection $P$ and final projection dominated by $\supp(c)$.
Taking the normal semifinite trace shows that
\[
 d_\tau((u-s)_+)=\tau(P)\le\tau(\supp(c))\le k.
\]
This finite bound uses the same $k$ for every $\tau$. Consequently,
if $d_\tau(v)=\infty$, then
$d_\tau((u-s)_+)<d_\tau(v)$ as well. For the remaining traces,
monotonicity gives
\[
 d_\tau((u-s)_+)\le d_\tau(u)<d_\tau(v).
\]
Definition~\ref{u:def:strict-comparison} now yields
$(u-s)_+\precsim v$ for every $s>0$. For each $n$, choose a
comparison witness $y_n$ such that
\[
 \norm{y_n^*vy_n-(u-1/n)_+}<1/n.
\]
Since $\norm{u-(u-1/n)_+}\le1/n$, these witnesses prove
$u\precsim v$.
\end{proof}

\subsection{The uniform tracial ultrapower}

Fix a free ultrafilter $\omega$ on $\N$. For a tracial state
$\tau$ on any unital $C^*$-algebra, and a nonempty set $X$ of such
states, write
\[
 \norm{x}_{2,\tau}=\tau(x^*x)^{1/2},\qquad
 \norm{x}_{2,X}=\sup_{\tau\in X}\norm{x}_{2,\tau}.
\]
Define
\[
 I=\left\{(a_l)\in\ell^\infty(\N,A):
             \lim_{l\to\omega}\norm{a_l}_{2,T}=0\right\},
 \qquad D=\ell^\infty(\N,A)/I.
\]
The inequalities
\[
 \norm{xy}_{2,\tau}\le\norm{x}\norm{y}_{2,\tau},
 \qquad
 \norm{yx}_{2,\tau}\le\norm{x}\norm{y}_{2,\tau},
 \qquad
 \norm{x^*}_{2,\tau}=\norm{x}_{2,\tau}
\]
show that $I$ is a two-sided $*$-ideal. It is norm closed because
$\norm{x}_{2,\tau}\le\norm{x}$. Thus $D$ is a unital $C^*$-algebra.
Faithfulness of the traces makes the map from $A$ given by constant
sequences injective; we henceforth identify $A$ with its image in $D$.

Let $\Lambda=T_\omega(A)$ be the set of limit traces on $D$. Its
members are the tracial states
\[
 (a_l)+I\longmapsto\lim_{l\to\omega}\tau_l(a_l),
 \qquad \tau_l\in T.
\]
Cauchy--Schwarz makes this formula independent of the representative,
and positivity, traciality and normalization follow coordinatewise.
For every bounded lift $(a_l)$ of $x\in D$ one has
\begin{equation}
\label{u:eq:limit-trace-norm}
 \begin{split}
 \norm{x}_{2,\Lambda}
   &=\lim_{l\to\omega}\norm{a_l}_{2,T},\\
 \sup_{\tau\in\Lambda}|\tau(x)|
   &=\lim_{l\to\omega}\sup_{\sigma\in T}|\sigma(a_l)|.
 \end{split}
\end{equation}
For either identity, the inequality from left to right follows
directly from the definition of a limit trace. For the reverse
inequality choose at each coordinate a trace whose value is within
$1/l$ of the corresponding supremum. For the first identity apply
this observation to $a_l^*a_l$ and take square roots; for the second
use continuity of absolute value under bounded scalar ultralimits.
In particular, $\norm{x}_{2,\Lambda}=0$ implies $x=0$ in $D$.
These identities concern the specified set $\Lambda$ and require
no description of the other traces on $D$.

Whenever $\tau\in\Lambda$ is evaluated on $D^{**}$, we mean its
canonical normal extension from $D$, again denoted by $\tau$.
It is tracial by the same separate continuity argument used above.
Neither this extension nor the extension of a faithful trace from
$A$ to $A^{**}$ is asserted to be faithful on the entire bidual.
All spectral trace inequalities below use positivity and normality
of the extension, which suffice for those inequalities.

\subsection{A bounded diagonal principle}

We give explicitly the form of the $\varepsilon$-test used below.
A $*$-polynomial with parameters in $D$ is a finite linear
combination of words in the variables, their adjoints, and finitely
many fixed elements of $D$ and their adjoints.

\begin{lemma}[Tracial diagonal principle]
\label{u:lem:tracial-diagonal}
Fix finitely many variables $x_1,\ldots,x_m$. For each variable
prescribe one of the following domains, with a fixed radius
$0\le R_i<\infty$: the closed ball of $D$, its selfadjoint part, or its positive
part. Let $(P_k)_{k\ge1}$ be countably many $*$-polynomials in these
variables with parameters in $D$. For each $k$ choose a test
\[
 t_k(x)=\norm{P_k(x)}_{2,\Lambda}
 \quad\text{or}\quad
 t_k(x)=\sup_{\tau\in\Lambda}|\tau(P_k(x))|,
\]
and a number $c_k\ge0$.
Suppose that for every $n\ge1$ and every $\varepsilon>0$ there is
an admissible tuple $x\in D^m$ with
\[
 t_k(x)<c_k+\varepsilon\qquad(1\le k\le n).
\]
Then there is an admissible tuple with $t_k(x)\le c_k$ for every
$k$. In particular, when every $c_k=0$, all the prescribed
tracial equations hold exactly. A polynomial with zero uniform
$2$-norm is zero as an element of $D$.
\end{lemma}

\begin{proof}
First, elements in each prescribed domain have lifts in the
corresponding coordinate domains with the same radius. For a
selfadjoint element, symmetrize an arbitrary lift and apply the
continuous clipping function to $[-R_i,R_i]$. For a positive element
clip to $[0,R_i]$ instead. For a general element with radius
$R_i>0$, if $a$ is any lift, replace it by $a f(a^*a)$, where
\[
 f(t)=\begin{cases}
 1,&0\le t\le R_i^2,\\
 R_i/\sqrt t,&t\ge R_i^2.
 \end{cases}
\]
This continuous functional calculus has norm at most $R_i$ and
does not change the quotient element. Radius zero has only the
zero lift.

Only countably many parameters occur in the polynomials. Fix one
bounded representative sequence for each of them, once and for
all. No multiplicative compatibility among these chosen lifts is
required. Given an admissible lifted tuple, let $t_{k,l}$ be the
coordinate test obtained by evaluating the lifted polynomial in
$A$ and replacing $\Lambda$ by $T$. Equation~\eqref{u:eq:limit-trace-norm}
and evaluation of polynomials in the quotient give
\[
 t_k(x)=\lim_{l\to\omega}t_{k,l}.
\]
Each of these coordinate sequences is bounded, since its
polynomial involves only finitely many bounded parameters and
variables in fixed balls.

For each $n$, choose an admissible trial tuple $x^{(n)}$ with
$t_k(x^{(n)})<c_k+1/(2n)$ for $k\le n$, and choose lifts
$x^{(n)}_{i,l}$ in the prescribed coordinate domains. The set
\[
 E_n=\left\{l:t_{k,l}(x^{(n)}_{1,l},\ldots,x^{(n)}_{m,l})
                     <c_k+1/n\text{ for all }k\le n\right\}
\]
belongs to $\omega$. Put
\[
 F_n=\{l\ge n\}\cap\bigcap_{j=1}^n E_j.
\]
These are decreasing members of $\omega$ with empty intersection.
For each $l$ that belongs to at least one $F_n$, let
$r(l)=\max\{n:l\in F_n\}$, which is finite and at most $l$,
and set $x_{i,l}=x^{(r(l))}_{i,l}$. At the remaining coordinates
set $x_{i,l}=0$.

The resulting sequences remain in their fixed balls. If $n\ge k$
and $l\in F_n$, then $r(l)\ge n$ and $l\in E_{r(l)}$, so
\[
 t_{k,l}(x_{1,l},\ldots,x_{m,l})
    <c_k+1/r(l)\le c_k+1/n.
\]
Taking the ultralimit, and then letting $n\to\infty$, proves the
claim. The last assertion follows from
Equation~\eqref{u:eq:limit-trace-norm}.
\end{proof}

The lemma also explains precisely where separability is used.
Conditions involving every element of a norm-separable algebra
can be imposed first on countable norm-dense subsets of the
appropriate balls. On fixed balls, polynomial evaluation is
uniformly norm continuous in each coefficient, and both test
seminorms are bounded by the operator norm. Thus the resulting
identities and inequalities extend to every coefficient by
continuity. For example,
\[
 \norm{[x,a]-[x,b]}_{2,\Lambda}
       \le2\norm{x}\norm{a-b},\qquad
 \sup_{\tau\in\Lambda}|\tau((a-b)x)|
       \le\norm{a-b}\norm{x}.
\]
In particular, a zero $2$-norm commutator gives exact commutation
in $D$. The number of variables and their norm bounds are fixed
throughout each application of Lemma~\ref{u:lem:tracial-diagonal}.

\subsection{The absorption criterion}

\begin{definition}[Uniform property $\Gamma$]
\label{u:def:uniform-gamma}
For a simple separable unital $A$ with $T(A)\ne\varnothing$, uniform
property $\Gamma$ means that for every $m\ge1$ there are
projections $p_1,\ldots,p_m\in D\cap A'$ satisfying
\[
 \sum_{i=1}^m p_i=1,
 \qquad
 \tau(ap_i)=\frac{1}{m}\tau(a)
 \quad(a\in A,\ \tau\in\Lambda,\ 1\le i\le m).
\]
This is \cite[Definition~2.1 and Equation~(2.1)]{CETW}, in the
uniform tracial ultrapower notation used here.
\end{definition}

\begin{theorem}[Castillejos--Evington--Tikuisis--White]
\label{u:thm:gamma-criterion}
Let $A$ be simple, separable, unital and nuclear, with
$T(A)\ne\varnothing$. If $A$ has strict comparison in the
functional formulation and uniform property $\Gamma$, then
\[
 A\cong A\otimes_{\min}\Zalg.
\]
\end{theorem}

This is the implication (ii)$\Rightarrow$(i) of
\cite[Theorem~5.6]{CETW}. Together with
Lemma~\ref{u:lem:comparison-convention}, it applies to the strict
comparison assumed in Theorem~\ref{u:thm:main}. We shall establish
uniform property $\Gamma$ in the following sections.
